% Standalone academic book / thesis template.
% For multi-chapter books driven from Markdown, prefer beagle-latex:quarto-book.
% This template is for hand-authored LaTeX books, slimmed from a full reference setup.

\documentclass[11pt,twoside,openright]{book}

% --- encoding & fonts ---
\usepackage[utf8]{inputenc}
\usepackage[T1]{fontenc}
\usepackage{newpxtext}
\usepackage{newpxmath}
\usepackage[scaled=0.95]{inconsolata}

% --- page layout ---
\usepackage[a4paper,
    inner=1in, outer=1.3in,
    top=1in, bottom=1.4in,
    bindingoffset=0.5in,
    headheight=14pt]{geometry}
\usepackage[final, protrusion=true, expansion=true]{microtype}
\usepackage{setspace}
\linespread{1.35}
\usepackage{emptypage}

% --- math, graphics, tables, lists ---
\usepackage{mathtools, amssymb, amsthm}
\usepackage{graphicx}
\usepackage[font=small, labelfont=bf, format=hang]{caption}
\usepackage{subcaption}
\usepackage{booktabs, array, multirow}
\usepackage{enumitem}
\setlist[itemize]{nosep, leftmargin=*, topsep=2pt, partopsep=0pt}
\setlist[enumerate]{nosep, leftmargin=*, topsep=2pt, partopsep=0pt}

\widowpenalty=10000
\clubpenalty=10000
\brokenpenalty=10000

% --- colors ---
\usepackage{xcolor}
\definecolor{chapterblue}{HTML}{1E3A5F}
\definecolor{linkblue}{RGB}{0,51,153}

% --- headers / footers ---
\usepackage{fancyhdr}
\pagestyle{fancy}
\fancyhf{}
\fancyhead[LE]{\small\slshape\nouppercase{\leftmark}}
\fancyhead[RO]{\small\slshape\nouppercase{\rightmark}}
\fancyfoot[C]{\small\thepage}
\renewcommand{\headrulewidth}{0.4pt}
\fancypagestyle{plain}{\fancyhf{}\fancyfoot[C]{\small\thepage}\renewcommand{\headrulewidth}{0pt}}

% --- chapter / section styling ---
\usepackage{titlesec}
\titleformat{\chapter}[display]
  {\normalfont\huge\bfseries\color{chapterblue}}
  {\chaptertitlename\ \thechapter}{20pt}{\Huge}
\titlespacing*{\chapter}{0pt}{-20pt}{40pt}

% --- theorems ---
\theoremstyle{plain}
\newtheorem{theorem}{Theorem}[chapter]
\newtheorem{lemma}[theorem]{Lemma}
\newtheorem{proposition}[theorem]{Proposition}
\newtheorem{corollary}[theorem]{Corollary}
\theoremstyle{definition}
\newtheorem{definition}[theorem]{Definition}
\newtheorem{example}[theorem]{Example}
\theoremstyle{remark}
\newtheorem{remark}[theorem]{Remark}

% --- algorithms ---
\usepackage{algorithm}
\usepackage{algpseudocode}

% --- TOC, index, citations, links ---
\usepackage{tocloft}
\setlength{\cftbeforechapskip}{1em}
\usepackage{imakeidx}
\makeindex[intoc]
\usepackage{natbib}
\usepackage{bookmark}
\usepackage{hyperref}
\hypersetup{
  colorlinks=true,
  linkcolor=chapterblue,
  citecolor=linkblue,
  urlcolor=linkblue,
  bookmarks=true, bookmarksnumbered=true, bookmarksopen=true,
}
\usepackage{cleveref}

% --- math shortcuts ---
\newcommand{\R}{\mathbb{R}}
\newcommand{\N}{\mathbb{N}}
\newcommand{\Z}{\mathbb{Z}}
\DeclareMathOperator*{\argmax}{arg\,max}
\DeclareMathOperator*{\argmin}{arg\,min}

% --- book metadata ---
\title{Book Title}
\author{Author Name}
\date{\the\year}

\begin{document}

% =========================================================
% FRONT MATTER (roman numerals, no chapter numbering)
% =========================================================
\frontmatter

\thispagestyle{empty}
\begin{center}
  \vspace*{1in}
  {\Huge\bfseries\color{chapterblue} Book Title} \\[0.5cm]
  {\LARGE Subtitle} \\[2cm]
  {\Large\itshape Author Name} \\[3cm]
  {\large Publisher} \\[0.3cm]
  {\large \the\year}
\end{center}
\cleardoublepage

\tableofcontents
\cleardoublepage

\chapter*{Preface}
\addcontentsline{toc}{chapter}{Preface}
Set the stage: who wrote this, why, and how it should be read.

\cleardoublepage

% =========================================================
% MAIN MATTER (arabic numerals, chapter numbering)
% =========================================================
\mainmatter

\chapter{Introduction}
\label{ch:introduction}

Open the book. Establish what the reader will get and how chapters are organized.

\section{Motivation}
Why does this topic matter?

\section{Scope}
What's in scope, what isn't, and how this differs from related books.

\chapter{Foundations}
\label{ch:foundations}

\begin{definition}[Core Concept]
\label{def:core}
A formal definition of the central object of study.
\end{definition}

\begin{theorem}[Main Result]
\label{thm:main}
A statement of the headline result.
\end{theorem}

\begin{proof}
Proof sketch or full proof.
\end{proof}

% Add additional chapters with \include{} or \input{}:
% \include{chapters/03-applications}
% \include{chapters/04-extensions}

% =========================================================
% BACK MATTER
% =========================================================
\backmatter

\appendix

\chapter{Mathematical Background}
\label{app:math}
Concepts assumed throughout the book.

% --- Bibliography ---
% Use a .bib file with natbib + bibtex (compile via beagle-latex:latex-bibliography):
\bibliographystyle{plainnat}
\bibliography{references}

% --- Index ---
\printindex

\end{document}
